{\large\textbf{Vortices in superfluid}}

\textbf{Introduction}

Superfluidity is a property of flowing without friction. Everyday experience
tells us that motion of an ordinary fluid (say, water at room temperature) is
always accompanied by viscous dissipation of energy, so that the flow gradually
becomes slower unless it is maintained by external forces. In contrast,
superfluid exhibits no loss of kinetic energy: once excited the motion of
superfluid can continue indefinitely. Superfluidity was originally discovered
experimentally in liquid helium.

We study properties of superfluid helium at zero temperature. It will be
treated as an incompressible fluid with density $\rho$. Flow continuity (the
fact that the mass flowing into and the mass flowing out of a given
infinitesimal volume are equal) implies that the flux of helium velocity
$\vec{v}$ through a closed surface is always zero. Superfluid velocity in this
aspect is analogous to the magnetic field intensity. By analogy with the
magnetic field lines, ``streamlines'' are tangential to the fluid velocity at
each point and their density is proportional to its magnitude.

True superflow has an important property of being irrotational: circulation of
superfluid velocity $\vec{v}$ along any closed path within helium is zero
\begin{equation}
\int\limits_L \vec{v} \cdot d\vec{l} = 0
\tag{1}
\end{equation}

This statement must be amended if superfluidity is absent along a thin
``vortex filament''. The thickness of the filament itself is of approximately
atomic dimensions $a$, but the vortex induces long range velocity field in
surrounding superfluid: velocity circulation around such filament is the
circulation quantum\footnote{Circulation quantization is a macroscopic quantum
effect and corresponds to the angular momentum quantization in Bohr model. The
circulation quantum can be expressed as $\kappa = \hbar/m_{\mathrm{He}}$, where
$m_{\mathrm{He}}$ is the mass of helium atom.}
\begin{equation}
\left| \int\limits_L \vec{v} \cdot d\vec{l} \right| = 2\pi\kappa,
\tag{2}
\end{equation}
and zero if the path can be contracted to a single point without crossing the
vortex, see Fig. 1. This supports the analogy between superflow and the
magnetic field created by wires carrying current: superposition of two valid
velocity distributions is a valid velocity distribution and the velocity at any
point is equal (up to a dimensional factor) to the magnetic field produced by
the unit currents running through a system of wires representing vortex
filaments.

\begin{center}\includegraphics[width=0.31\linewidth]{figures/APhO_2017_Q1_fig1.png}\end{center}

Fig. 1: Vortex filament (red) in superfluid (light blue). Velocity circulations
along paths $L_1$, $L_2$, $L_5$, and $L_6$ are all zero, but those for $L_3$
and $L_4$ are equal to $\pm 2\pi\kappa$. Note that circulations along $L_3$ and
$L_4$ have opposite signs.

\textbf{Part A. Steady filament (0.75 points)}

Consider a cylindrical beaker (radius $R_0 \gg a$) of superfluid helium and a
straight vertical vortex filament in its center Fig. 2.

\textbf{A.1}\quad Plot the streamlines. Find out the velocity $v$ at a point
$\vec{r}$. \hfill 0.25pt

\textbf{A.2}\quad Work out the free surface shape (height as a function of
coordinate $z(\vec{r})$) around the vortex. Free fall acceleration is $g$.
Surface tension can be neglected. \hfill 0.5pt

\begin{center}\includegraphics[width=0.31\linewidth]{figures/APhO_2017_Q1_fig2.png}\end{center}

\begin{center}Fig. 2: Straight vortex along the axis of a beaker.\end{center}

\textbf{Part B. Vortex motion (1.4 points)}

Free vortices move about in space with the flow\footnote{This is a consequence
of momentum conservation, see next section.}. In other words each element of
the filament moves with the velocity $\vec{v}$ of the fluid at the position of
that element.

As an example, consider a pair of counter-rotating straight vortices placed
initially at distance $r_0$ from each other, see Fig. 3. Each vortex produces
velocity $v_0 = \kappa/r_0$ at the axis of another. As a result, the vortex
pair moves rectilinearly with constant speed $v_0 = \kappa/r_0$ so that the
distance between them remains unchanged.

\begin{center}\includegraphics[width=0.31\linewidth]{figures/APhO_2017_Q1_fig3.png}\end{center}

\begin{center}Fig. 3: Parallel vortex filaments with opposite circulations.\end{center}

\textbf{B.1}\quad Consider two identical straight vortices initially placed at
distance $r_0$ from each other as shown in Fig. 4. Find initial velocities of
the vortices and draw their trajectories. \hfill 0.25pt

\begin{center}\includegraphics[width=0.31\linewidth]{figures/APhO_2017_Q1_fig4.png}\end{center}

\begin{center}Fig. 4: Parallel vortex filaments with equal circulations.\end{center}

A beaker of helium (see Part A) is filled with triangular lattice
($u \ll R_0$) of identical vertical vortices, see Fig. 5.

\begin{center}\includegraphics[width=0.51\linewidth]{figures/APhO_2017_Q1_fig5.png}\end{center}

\begin{center}Fig. 5: Triangular lattice of vortices in a beaker. The view from above.\end{center}

\textbf{B.2}\quad Draw the trajectories of vortices A, B, and C (located in the
center). \hfill 0.15pt

\textbf{B.3}\quad Find velocity $v(\vec{r})$ of a vortex positioned at
$\vec{r}$. \hfill 0.4pt

\textbf{B.4}\quad Find the distance AB$(t)$ between the vortices A and B at
time $t$. Treat AB$(0)$ as given. \hfill 0.35pt

\textbf{B.5}\quad Work out the \textquotedbl smoothed out\textquotedbl{} (omitting the lattice structure)
free helium surface shape $z(\vec{r})$. \hfill 0.25pt

\textbf{Part C. Momentum and energy (1.75 points)}

The long range velocity field is the major contribution to the energy of a
system of vortices, it is insensitive to exact structure of the filament. The
filament itself can not be properly described by the macroscopic theory and
apparent singularities (infinities) are insignificant. Real physical
quantities, such as the energy, of the region inside a thin tube of radius $a$
around the filament should be neglected. Outside of this tube the density of
superflow kinetic energy $\rho v^2/2$ (where $\rho = \mathrm{const}$) is
analogous to the energy density of the magnetic field $B^2/(2\mu_0)$ --- they
are both quadratic in respective variables. This analogy together with the
correspondence between magnetic field and superfluid velocity generated by
vortices (currents) facilitates calculation of the flow energy for a given
system. For instance, given the inductance of a circular wire loop
$L \approx \mu_0 R \log(R/a)$, where $R$ is the loop radius and $a$ is wire
radius, we get the superfluid vortex loop energy\footnote{This expression is
also valid only if $\log R/a \gg 1$.}
\begin{equation}
U \approx 2R\rho\pi^2\kappa^2 \log(R/a)
\tag{3}
\end{equation}

Total fluid momentum is also determined by the long range velocity
distribution. It is obtained by integration of the momentum density
$\rho\vec{v}$. Again, consider a flow generated by a circular vortex loop placed
in $xy$ plane. It is obvious from the symmetry considerations, that total
momentum has only $z$ component:
\begin{equation}
P = \int \rho v_z dV = \rho \int\int \underbrace{\left(\int v_z dz\right)}_{q(x,y)} dxdy
\tag{4}
\end{equation}

The innermost integration is in fact an integration along appropriate paths
parallel to $z$-axis, see Fig. 6. From the circulation identity (2) it follows
that
\begin{equation}
q(x,y) = \int\limits_{L(x,y)} \vec{v} \cdot d\vec{l}
\tag{5}
\end{equation}
is piecewise constant. Particularly, it is zero for paths passing outside the
ring and $2\pi\kappa$ for paths inside it. Total momentum is therefore
\begin{equation}
P = \rho \cdot \pi R^2 \cdot 2\pi\kappa = 2\pi^2\rho R^2\kappa
\tag{6}
\end{equation}

\begin{center}\includegraphics[width=0.92\linewidth]{figures/APhO_2017_Q1_fig6.png}\end{center}

Fig. 6: Velocity field of a circular vortex loop and integration paths (green)
for $q(x,y)$ calculation.

\begin{center}\includegraphics[width=0.41\linewidth]{figures/APhO_2017_Q1_fig7.png}\end{center}

\begin{center}Fig. 7: A nearly rectangular vortex loop, $b \ll d$.\end{center}

\textbf{C.1}\quad Consider a nearly rectangular vortex loop $b \times d$,
$b \ll d$, Fig. 7. Indicate the direction of its momentum $\vec{P}$. Find out
the momentum magnitude. \hfill 0.3pt

\textbf{C.2}\quad Calculate its energy $U$. \hfill 0.7pt

\textbf{C.3}\quad Suppose we shift a long straight vortex filament by a distance
$b$ in $x$ direction, see Fig. 8. How much does the fluid momentum change?
Indicate the momentum change direction. The filament length (constrained by the
vessel walls) is $d$. \hfill 0.75pt

\begin{center}\includegraphics[width=0.30\linewidth]{figures/APhO_2017_Q1_fig8.png}\end{center}

\begin{center}Fig. 8: Momentum changes whenever the vortex shifts with respect to the fluid.\end{center}

\textbf{Part D. Trapped charges (2.85 points)}

Electrons, if injected in helium, get trapped in the vortex filaments. Here
and below polarizability of helium can be neglected ($\epsilon = 1$).

\begin{center}\includegraphics[width=0.51\linewidth]{figures/APhO_2017_Q1_fig9.png}\end{center}

\begin{center}Fig. 9: Straight vortex in a uniform electric field.\end{center}

\textbf{D.1}\quad Consider a straight vortex charged with uniform linear
density $\lambda < 0$ in a uniform electric field $\vec{E}$. Draw the vortex
trajectory. Find its velocity as a function of time. \hfill 0.5pt

A circular vortex loop of radius $R_0$ initially charged with uniform linear
density $\lambda < 0$ is placed in a uniform electric field $\vec{E}$
perpendicular to its plane, opposite to its momentum $\vec{P}_0$.

\begin{center}\includegraphics[width=0.72\linewidth]{figures/APhO_2017_Q1_fig10.png}\end{center}

\begin{center}Figure 10: (left) Vortex ring in a uniform electric field. (right) Cross section of the ring.\end{center}

\textbf{D.2}\quad Draw the trajectory of the loop center $C$. Find the radius
of the loop as a function of time. \hfill 0.6pt

\textbf{D.3}\quad Find its velocity $v(t)$ as a function of time. \hfill 1.5pt

\textbf{D.4}\quad The field is switched off at a time $t^*$ when the velocity
reaches the value $v^* = v(t^*)$. Find the loop velocity $v(t)$ at a later
time $t > t^*$. \hfill 0.25pt

\textbf{Part E. Influence of the boundaries (3.25 points)}

Solid walls alter the velocity field created by a vortex filament, because the
fluid cannot flow through them. Mathematically this means that the wall-normal
velocity component vanishes at the wall surface.

\begin{center}\includegraphics[width=0.31\linewidth]{figures/APhO_2017_Q1_fig11.png}\end{center}

\begin{center}Fig. 11: Straight vortex filament near a flat wall.\end{center}

\textbf{E.1}\quad Draw the trajectory of a straight vortex, initially placed at
a distance $h_0$ from a flat wall. Find its velocity as a function of time.
\hfill 0.5pt

Consider a straight vortex placed in a corner at a distance $h_0$ from both
walls.

\begin{center}\includegraphics[width=0.31\linewidth]{figures/APhO_2017_Q1_fig12.png}\end{center}

\begin{center}Fig. 12: Straight vortex filament in a corner.\end{center}

\textbf{E.2}\quad What is the initial velocity $v_0$ of the vortex? \hfill 0.75pt

\textbf{E.3}\quad Draw the trajectory of the vortex. \hfill 0.5pt

\textbf{E.4}\quad What is the velocity of the vortex $v_\infty$ after very long
time? \hfill 1.5pt
